\chapter{Introduction}
\label{ch:introduction}

\section{Background and Motivation}
Biological research in the modern era relies increasingly on multi-omics profiling and computational data modeling. In this chapter, we outline the foundational principles and state-of-the-art literature motivating this dissertation \cite{Horvath2013}.

\section{Heading Hierarchy Example}
As required by the institutional formatting rubric (Sec III.J), headings must strictly adhere to the four-level hierarchy:
\begin{enumerate}
  \item \textbf{Primary Headings} (\verb|\chapter{}|): ALL CAPITAL LETTERS, bold, centered, double spaced above major sections.
  \item \textbf{Secondary Headings} (\verb|\section{}|): Sentence Case, bold, left-aligned.
  \item \textbf{Tertiary Headings} (\verb|\subsection{}|): Sentence Case, italicized, unbolded, left-aligned.
  \item \textbf{Quaternary Headings} (\verb|\subsubsection{}|): Sentence Case, italicized, in-line with paragraph, followed by a period.
\end{enumerate}

\subsection{Tertiary Subsection Heading Example}
Tertiary headings are italicized and unbolded in sentence case. Paragraph text begins double-spaced underneath.

\subsubsection{Quaternary Heading Example}
Quaternary headings run directly in line with paragraph text, followed immediately by a period.

\section{Table Formatting Example}
Per Rubric Sec III.A, table titles/captions must be placed \textbf{above} the table, single-spaced, with bolded table references. Cohort characteristics are summarized in \TABLE{demographics}.

\begin{table}[htb]
  \centering
  \caption[Sample demographic summary table.]{Demographic summary of study cohorts evaluated across chapters.}
  \label{tab:demographics}
  \begin{tabular}{lccc}
    \toprule
    \textbf{Cohort Name} & \textbf{Ancestry} & \textbf{Sample Size (N)} & \textbf{Platform} \\
    \midrule
    Discovery Cohort     & Multi-Ancestry    & 621                      & EPIC BeadChip \\
    Replication Cohort A & African American  & 1,394                    & 450K BeadChip \\
    Replication Cohort B & European          & 729                      & 450K BeadChip \\
    \bottomrule
  \end{tabular}
\end{table}

\section{Figure Formatting Example}
Per Rubric Sec III.B, figure titles and legends must be placed \textbf{below} the figure and single-spaced. The analytical workflow is outlined in \FIG{workflow}, and its preprocessing step in \FIG[A]{workflow}.

\begin{figure}[htb]
  \centering
  % Replace with your figure graphic (e.g., \includegraphics[width=0.85\linewidth]{figures/fig1.pdf}),
  % and describe it in figure_alt_text.tex: \FigAlt{fig1}{<what it shows>}
  \rule{0.8\linewidth}{2in} % Placeholder graphic box
  \caption[Schematic overview of the analytical study workflow.]{
    \textbf{Schematic overview of the analytical study workflow.}
    High-throughput DNA methylation and genomic variant data are preprocessed, followed by ancestry-stratified evaluation of epigenetic clock predictions and methylation quantitative trait loci (meQTL) mapping.
  }
  \label{fig:workflow}
\end{figure}

\section{Figure and Table Reference Example}
Every in-text reference to a figure or table is bold, panel letter included (Rubric Sec III.A.7). Write references with the macros, never by hand:
\begin{itemize}
  \item \verb|\FIG{workflow}| gives \FIG{workflow}; \verb|\FIG[A]{workflow}| gives \FIG[A]{workflow}.
  \item \verb|\TABLE{demographics}| gives \TABLE{demographics}.
  \item \verb|\cref{fig:workflow}| gives \cref{fig:workflow} (bolded by the preamble's \verb|\crefformat| block).
\end{itemize}
Floats use \verb|[htb]| (or no option). Never \verb|[p]| or \verb|[htbp]|: a float page in the middle of running text is not allowed.

\section{Block Quotation Formatting Example}
Per Rubric Sec I.F and III.C, prose quotations longer than four lines must be single-spaced and indented uniformly from the left margin, as in the following example:

\begin{quotation}
  Biological aging is a complex, multifaceted process characterized by the gradual deterioration of physiological function and increased susceptibility to disease. Epigenetic modifications, particularly DNA methylation at cytosine-phosphate-guanine dinucleotides, provide a dynamic and quantifiable record of cellular chronological and physiological history.
\end{quotation}

Double-spaced prose continues here immediately following the quotation block.
